\chapter{凸泛函与无穷维微积分}

\label{chap:05}

第 \ref{chap:04} 章把凸集、分离与支撑超平面的几何语言建立起来之后，下一步自然要把视角从“集合”转向“泛函”。在有限维凸优化里，这一步往往表现为熟悉的函数图像、梯度、次梯度和最优性条件；但在无穷维中，真正困难的并不是把记号从 $x\in\mathbb R^n$ 改成 $u\in X$，而是要重新界定哪些连续性足以支撑存在性，哪些可微性足以支撑一阶展开，以及哪些局部线性信息能够在不可微情形下留下来。

本章正是为这些问题建立局部语言。第 \ref{sec:5-1} 节先讨论下半连续性、弱紧性与直接法，说明凸泛函的极小值存在绝不是一个无条件口号，而必须由拓扑、强制性和紧性来源共同支撑。第 \ref{sec:5-2} 节再把方向导数、Gâteaux 微分与 Fréchet 微分区分开来，解释为什么无穷维里“可微”本身就有层级。第 \ref{sec:5-3} 节引入对偶空间中的次微分，把不可微凸泛函的局部信息重写成支撑超平面的函数化版本。最后第 \ref{sec:5-4} 节整理次微分演算，给出和规则、线性映射下的拉回以及 Moreau--Rockafellar 和公式，为第 \ref{chap:06} 章的共轭理论和第 9 至 10 章的单调算子与分裂算法提供可直接调用的接口。

阅读本章时，需要始终记住三层环境的区别。局部凸或 Banach 框架已经足够谈凸泛函、连续线性泛函与分离；自反 Banach 空间开始提供弱紧性，因此能承载直接法；Hilbert 空间则通过第 \ref{sec:3-1} 节定理~\ref{thm:riesz-representation-hilbert} 把导数和对偶重新识别为向量，使梯度、投影、邻近映射与算法迭代获得更具体的几何解释。后文关于 Fenchel--Rockafellar 对偶、KKT 条件、极大单调算子与 ADMM 的很多论证，都会回链到这一章。

\section{下半连续性与极值定理}

\label{sec:5-1}

若把凸优化理解为“在一个几何良好的空间里最小化一个几何良好的目标”，那么存在性问题就是第一道关口。有限维情形中，闭有界集紧、连续函数在紧集上取极值，这一链条几乎不必单独说明；但在无穷维里，闭有界集通常不紧，甚至极小化列都可能完全逃逸。因此，本节要把“极小值存在”重写成真正可复用的标准模板：目标泛函需要适当的下半连续性，可行区域需要某种弱紧性来源，而强制性则负责阻止极小化过程向无穷远逃走。

\subsection{Proper 凸泛函与下半连续性}

\begin{definition}[Proper 凸泛函]\label{def:proper-functional}
设 $X$ 为实向量空间，函数
\[
f\colon X\to \mathbb R\cup\{+\infty\}
\]
称为\emph{proper} 的，若其有效定义域
\[
\operatorname{dom}f:=\{x\in X:f(x)<+\infty\}
\]
非空，且 $f$ 不取值 $-\infty$。若再满足对任意 $x,y\in X$ 与任意 $\lambda\in[0,1]$，
\[
f(\lambda x+(1-\lambda)y)\le \lambda f(x)+(1-\lambda)f(y),
\]
则称 $f$ 为一个\emph{proper 凸泛函}。
\end{definition}

\begin{definition}[下半连续与弱下半连续]\label{def:lower-semicontinuity-functional}
设 $X$ 为拓扑向量空间，$f\colon X\to \mathbb R\cup\{+\infty\}$ 为一个 proper 泛函。若对每个 $r\in\mathbb R$，次水平集
\[
\{x\in X:f(x)\le r\}
\]
都是闭集，则称 $f$ 在给定拓扑下\emph{下半连续}。当 $X$ 为 Banach 空间且所用拓扑是弱拓扑 $\sigma(X,X^\ast)$ 时，称 $f$ 为\emph{弱下半连续}。
\end{definition}

\begin{remark}
定义~\ref{def:lower-semicontinuity-functional} 与第 \ref{sec:1-1} 节命题~\ref{prop:closure-net} 等价地说明：$f$ 在某个拓扑下下半连续，当且仅当对任意收敛网 $x_\alpha\to x$ 都有
\[
f(x)\le \liminf_\alpha f(x_\alpha).
\]
后文一旦工作在弱拓扑或弱$^\ast$拓扑中，就必须把这条判别理解成关于网的陈述，而不能默认退回序列语言。
\end{remark}

\begin{definition}[强制性]\label{def:coercive-functional}
设 $X$ 为赋范空间。proper 泛函 $f\colon X\to \mathbb R\cup\{+\infty\}$ 称为\emph{强制}的，若
\[
\|x\|\to\infty \quad\Longrightarrow\quad f(x)\to+\infty.
\]
等价地说，对任意实数 $r$，次水平集
\[
\{x\in X:f(x)\le r\}
\]
都是有界的。
\end{definition}

\begin{remark}
强制性并不等于有极小值；它只保证极小化过程不会逃向无穷远。真正的存在性还需要结合弱紧性或紧嵌入等额外结构。第 \ref{sec:3-2} 节推论~\ref{cor:hilbert-weak-compactness} 已经说明：在 Hilbert 或更一般自反 Banach 空间中，闭球的弱紧性才是直接法得以启动的关键。
\end{remark}

\subsection{弱紧集上的极小值存在}

\begin{theorem}[弱紧集上的弱下半连续泛函取极小值]\label{thm:weak-lsc-attainment}
设 $X$ 为 Banach 空间，$C\subset X$ 为非空弱紧集，$f\colon X\to \mathbb R\cup\{+\infty\}$ 为 proper 且在 $C$ 上弱下半连续的泛函。则存在 $\bar x\in C$ 使
\[
f(\bar x)=\inf_{x\in C}f(x).
\]
\end{theorem}

\begin{proof}
记
\[
m:=\inf_{x\in C}f(x).
\]
由 proper 性可知 $m<+\infty$。对每个 $\varepsilon>0$ 选取 $x_\varepsilon\in C$ 使
\[
f(x_\varepsilon)\le m+\varepsilon.
\]
将 $(0,+\infty)$ 赋予反序，即 $\varepsilon_1\preceq \varepsilon_2$ 当且仅当 $\varepsilon_1\ge \varepsilon_2$，于是 $(x_\varepsilon)$ 构成一条以“精度趋于零”为指标的网。由于 $C$ 弱紧，它存在弱收敛子网，仍记为
\[
x_\alpha \rightharpoonup \bar x\in C.
\]
由定义~\ref{def:lower-semicontinuity-functional} 的网刻画，
\[
f(\bar x)\le \liminf_\alpha f(x_\alpha)\le \liminf_\alpha (m+\alpha)=m.
\]
另一方面 $m$ 本就是 $C$ 上的下确界，因此 $m\le f(\bar x)$。综上
\[
f(\bar x)=m,
\]
即 $\bar x$ 是所求极小元。
\end{proof}

\subsection{自反空间中的直接法}

\begin{theorem}[自反 Banach 空间中的直接法]\label{thm:direct-method-reflexive}
设 $X$ 为自反 Banach 空间，$f\colon X\to \mathbb R\cup\{+\infty\}$ 为 proper、弱下半连续且强制的凸泛函。则 $f$ 在 $X$ 上取到极小值。
\end{theorem}

\begin{proof}
记
\[
m:=\inf_{x\in X} f(x).
\]
由 proper 性可知 $m<+\infty$。同样对每个 $\varepsilon>0$ 选取 $x_\varepsilon\in X$ 满足
\[
f(x_\varepsilon)\le m+\varepsilon.
\]
由于 $f$ 强制，这条极小化网必有界；否则若 $\|x_\varepsilon\|\to\infty$，则 $f(x_\varepsilon)\to+\infty$，与上式矛盾。故存在 $R>0$ 使
\[
x_\varepsilon\in \overline{B}(0,R),\qquad \forall \varepsilon.
\]
由第 \ref{sec:3-2} 节定义~\ref{def:reflexive-space} 及 Banach--Alaoglu 定理的推论，自反 Banach 空间中的闭球在弱拓扑下紧，因此 $\overline{B}(0,R)$ 弱紧。由定理~\ref{thm:weak-lsc-attainment}，$(x_\varepsilon)$ 的某个子网弱收敛到某个 $\bar x\in \overline{B}(0,R)$，且
\[
f(\bar x)\le \liminf_\alpha f(x_\alpha)\le m.
\]
由 $m$ 的定义，又有 $m\le f(\bar x)$。故
\[
f(\bar x)=m.
\]
于是 $f$ 在 $X$ 上取到极小值。
\end{proof}

\begin{remark}
定理~\ref{thm:direct-method-reflexive} 是本书后文最常调用的存在性模板之一。正确的口号不是“下半连续凸泛函必有极小值”，而是：“若极小化过程被强制性压进某个弱紧区域，并且目标在该拓扑下下半连续，则直接法成立。”这一区分在第 \ref{sec:6-4} 节对偶间隙分析和第 \ref{sec:11-1} 节控制问题里都会反复出现。
\end{remark}

\subsection{典型例子}

\begin{example}[Sobolev 型能量泛函]\label{ex:sobolev-energy-functional}
在 Hilbert 空间 $H_0^1(\Omega)$ 上考虑
\[
J(u)=\frac12\int_\Omega |\nabla u(x)|^2\,dx+\int_\Omega \Phi(u(x))\,dx-\langle \ell,u\rangle,
\]
其中 $\Phi$ 为凸下半连续函数，$\ell\in H^{-1}(\Omega)$。把这个例子放在这里，是为了说明偏微分方程中的变分问题通常正是定理~\ref{thm:direct-method-reflexive} 的标准应用场景：二次梯度项提供强制性，凸势能与线性项决定弱下半连续性和可解性。
\end{example}

\begin{example}[有限维闭有界水平集]\label{ex:finite-dimensional-level-set}
当 $X=\mathbb R^n$ 时，弱拓扑与通常拓扑一致，闭有界集自动紧。于是若 $f$ 的某个次水平集
\[
\{x\in\mathbb R^n:f(x)\le r\}
\]
非空且闭有界，则极小值存在直接退化为 Weierstrass 定理。把这个例子放在这里，是为了说明 Boyd 书里许多“函数连续、水平集闭有界，因此解存在”的论证，实际上是定理~\ref{thm:weak-lsc-attainment} 在欧氏空间中的坍缩。
\end{example}

\begin{example}[Tikhonov 型正则化]\label{ex:tikhonov-regularization}
设 $H$ 为 Hilbert 空间，$A\colon H\to Y$ 为有界线性算子，考虑
\[
f(u)=\frac12\|Au-b\|_Y^2+\lambda \|u\|_H^2,\qquad \lambda>0.
\]
该泛函是 proper、凸、连续且强制的，因此由定理~\ref{thm:direct-method-reflexive} 在 $H$ 上存在极小元。把这个例子放在这里，是为了连接后文的逆问题、机器学习与第 \ref{sec:10-1} 节前向--后向分裂：正则化项不仅改善数值稳定性，也提供存在性所需的强制性。
\end{example}

\subsection*{有限维对照}

Boyd 第 3 章里关于凸函数、次水平集和极小值存在的叙述，默认都工作在 $\mathbb R^n$ 中，因此“闭有界即紧”这一事实被自动吞掉。本节的作用是把这一步显式拆开：定理~\ref{thm:weak-lsc-attainment} 负责把极小值存在写成“紧性 + 下半连续”的拓扑命题，定理~\ref{thm:direct-method-reflexive} 则说明在无穷维里，这个紧性往往要通过弱拓扑、自反性和强制性重新获得。

\subsection*{习题（待补）}

\begin{exercise}[弱下半连续判别占位]\label{exr:weak-lsc-placeholder}
补一题：要求把定义~\ref{def:lower-semicontinuity-functional} 用网重写，并结合命题~\ref{prop:closure-net} 说明为什么弱下半连续性的正确语言不是“子列极限”而是“子网极限”。
\end{exercise}

\begin{exercise}[直接法桥接题占位]\label{exr:direct-method-reflexive-placeholder}
补一题：在 Hilbert 空间中证明例~\ref{ex:tikhonov-regularization} 的极小元存在，并明确指出证明中分别在哪一步使用了强制性、推论~\ref{cor:hilbert-weak-compactness} 与下半连续性。
\end{exercise}


\section{方向导数：Gâteaux 微分与 Fréchet 微分}

\label{sec:5-2}

存在性解决以后，下一步自然是研究泛函在局部如何变化。有限维里，梯度既是最速上升方向，也是线性近似的系数，因此“可微”通常只剩一种标准形态；但在无穷维中，方向上的极限存在并不自动意味着存在一个统一的连续线性逼近。也就是说，局部一阶信息本身有强弱层级。第 \ref{sec:5-3} 节的次微分正是建立在这种区分之上，而第 10 章的一阶算法则更加依赖其中较强的一层。

\subsection{方向导数与 Gâteaux 微分}

\begin{definition}[方向导数]\label{def:directional-derivative}
设 $X$ 为实向量空间，$f\colon X\to \mathbb R\cup\{+\infty\}$，$x\in \operatorname{dom}f$，$h\in X$。若极限
\[
f'(x;h):=\lim_{t\downarrow 0}\frac{f(x+th)-f(x)}{t}
\]
存在于 $\mathbb R\cup\{\pm\infty\}$，则称其为 $f$ 在点 $x$ 沿方向 $h$ 的\emph{方向导数}。
\end{definition}

\begin{definition}[Gâteaux 微分]\label{def:gateaux-derivative}
设 $X,Y$ 为 Banach 空间，$U\subset X$ 为开集，$f\colon U\to Y$。若对每个 $h\in X$，极限
\[
\lim_{t\to 0}\frac{f(x+th)-f(x)}{t}
\]
存在，且存在连续线性算子 $A\in \mathcal L(X,Y)$ 使得
\[
\lim_{t\to 0}\frac{f(x+th)-f(x)}{t}=Ah,\qquad \forall h\in X,
\]
则称 $f$ 在 $x$ 处\emph{Gâteaux 可微}，并称 $A$ 为 $f$ 在 $x$ 处的 Gâteaux 导数，记为 $D_Gf(x)$。
\end{definition}

\begin{remark}
定义~\ref{def:gateaux-derivative} 中最微妙的一点在于：极限是对每个方向逐个成立的。它要求存在一个共同的连续线性算子来记录所有方向的导数，但并不要求余项对所有小扰动同时一致控制。这也是 Gâteaux 微分弱于 Fréchet 微分的根本原因。
\end{remark}

\subsection{Fréchet 微分与 Hilbert 梯度}

\begin{definition}[Fréchet 微分]\label{def:frechet-derivative}
设 $X,Y$ 为 Banach 空间，$U\subset X$ 为开集，$f\colon U\to Y$。若存在连续线性算子 $A\in\mathcal L(X,Y)$ 使
\[
\lim_{\|h\|\to 0}\frac{\|f(x+h)-f(x)-Ah\|_Y}{\|h\|_X}=0,
\]
则称 $f$ 在 $x$ 处\emph{Fréchet 可微}，并称 $A$ 为 $f$ 在 $x$ 处的 Fréchet 导数，记为 $Df(x)$。
\end{definition}

\begin{proposition}[Fréchet 可微蕴含 Gâteaux 可微]\label{prop:frechet-implies-gateaux}
设 $U\subset X$ 为 Banach 空间中的开集，$f\colon U\to Y$。若 $f$ 在 $x\in U$ 处 Fréchet 可微，且导数为 $Df(x)$，则 $f$ 在 $x$ 处 Gâteaux 可微，且
\[
D_G f(x)h=Df(x)h,\qquad \forall h\in X.
\]
\end{proposition}

\begin{proof}
由定义~\ref{def:frechet-derivative}，对任意 $h\in X$ 与任意标量 $t\neq 0$，
\[
f(x+th)-f(x)-Df(x)(th)=r(th),
\]
其中
\[
\frac{\|r(th)\|_Y}{|t|\|h\|_X}\to 0\qquad (t\to 0).
\]
故
\[
\frac{f(x+th)-f(x)}{t}=Df(x)h+\frac{r(th)}{t}.
\]
当 $t\to 0$ 时，右侧第二项趋于 $0$，于是
\[
\lim_{t\to 0}\frac{f(x+th)-f(x)}{t}=Df(x)h.
\]
这正是定义~\ref{def:gateaux-derivative}。
\end{proof}

\begin{remark}
命题~\ref{prop:frechet-implies-gateaux} 的逆命题一般不成立。也就是说，逐方向的导数存在并拼出一个线性算子，并不能保证余项在小球上受一致控制。对优化而言，这个差别非常重要：Gâteaux 可微往往足以给出形式上的一阶展开，但 Fréchet 可微才更接近数值算法所依赖的稳定线性近似。
\end{remark}

\begin{definition}[Hilbert 空间中的梯度]\label{def:hilbert-gradient}
设 $H$ 为 Hilbert 空间，$f\colon H\to \mathbb R$ 在 $x\in H$ 处 Fréchet 可微。由第 \ref{sec:3-1} 节定理~\ref{thm:riesz-representation-hilbert}，存在唯一向量 $\nabla f(x)\in H$ 使
\[
Df(x)h=\langle \nabla f(x),h\rangle_H,\qquad \forall h\in H.
\]
该向量称为 $f$ 在 $x$ 处的\emph{梯度}。
\end{definition}

\subsection{凸可微泛函的一阶信息}

\begin{proposition}[凸可微泛函的一阶支撑不等式]
设 $X$ 为 Banach 空间，$U\subset X$ 为凸开集，$f\colon U\to \mathbb R$ 为凸函数。若 $f$ 在 $x\in U$ 处 Gâteaux 可微，则对任意 $y\in U$ 都有
\[
f(y)\ge f(x)+D_Gf(x)(y-x).
\]
\end{proposition}

\begin{proof}
固定 $y\in U$，考虑一维函数
\[
\phi(t)=f(x+t(y-x)),\qquad t\in[0,1].
\]
由凸性知 $\phi$ 是区间上的凸函数。对任意 $t\in(0,1]$，凸函数的割线斜率单调，故
\[
\frac{\phi(t)-\phi(0)}{t}\ge \phi'_+(0).
\]
令 $t=1$ 并利用 Gâteaux 导数存在，得到
\[
f(y)-f(x)\ge D_Gf(x)(y-x).
\]
\end{proof}

\begin{remark}
这条支撑不等式说明：一旦凸泛函在某点可微，其导数就自动给出该点处的支撑超平面。第 \ref{sec:5-3} 节中的次微分正是要把这种“单值支撑”推广到不可微情形。
\end{remark}

\subsection{典型例子}

\begin{example}[积分型泛函的 Gâteaux 导数]\label{ex:gateaux-integral-functional}
设 $\Omega$ 为有限测度空间，$X=L^2(\Omega)$，并考虑
\[
f(u)=\int_\Omega \Phi(u(x))\,dx,
\]
其中 $\Phi\in C^1(\mathbb R)$ 且 $\Phi'$ 具有适当增长。则在合适的定义域上，
\[
D_G f(u)h=\int_\Omega \Phi'(u(x))h(x)\,dx.
\]
把这个例子放在这里，是为了说明 Gâteaux 导数自然落在对偶配对之中：方向导数看似是“沿 $h$ 的一维变化”，但最后得到的却是一个连续线性泛函。
\end{example}

\begin{example}[欧氏二次型的梯度]\label{ex:quadratic-gradient}
在 $\mathbb R^n$ 上令
\[
f(x)=\frac12 x^\top Qx+b^\top x+c,
\]
其中 $Q$ 为对称矩阵。则 $f$ Fréchet 可微且
\[
\nabla f(x)=Qx+b.
\]
把这个例子放在这里，是为了表明 Boyd 书里最熟悉的梯度公式，其实正是定义~\ref{def:hilbert-gradient} 在有限维 Hilbert 空间中的退化。
\end{example}

\begin{example}[Hilbert 空间中的最小二乘能量]\label{ex:hilbert-least-squares}
设 $H_1,H_2$ 为 Hilbert 空间，$A\colon H_1\to H_2$ 为有界线性算子。对
\[
f(u)=\frac12\|Au-b\|_{H_2}^2
\]
有
\[
Df(u)h=\langle Au-b,Ah\rangle_{H_2}
=\langle A^\ast(Au-b),h\rangle_{H_1},
\]
其中 $A^\ast$ 由第 \ref{sec:3-3} 节定义~\ref{def:adjoint-operator} 给出。于是
\[
\nabla f(u)=A^\ast(Au-b).
\]
把这个例子放在这里，是为了连接第 \ref{sec:3-3} 节的伴随算子与本节的梯度概念，也为后文最速下降和正则化算法提供标准原型。
\end{example}

\subsection*{有限维对照}

Boyd 第 9 章里把梯度、Jacobian 和最速下降法建立在欧氏空间几何之上，因此“可微”几乎默认就是 Fréchet 可微，而且导数自动由向量表示。本节的作用，是把这种有限维母语拆成三个层级：定义~\ref{def:directional-derivative} 只记录单方向变化，定义~\ref{def:gateaux-derivative} 记录逐方向的线性化，而定义~\ref{def:frechet-derivative} 与定义~\ref{def:hilbert-gradient} 则给出算法上真正稳定的一阶近似。

\subsection*{习题（待补）}

\begin{exercise}[Gâteaux 与 Fréchet 区分占位]\label{exr:gateaux-frechet-placeholder}
补一题：构造一个 Gâteaux 可微但并非 Fréchet 可微的泛函例子，并说明这种失效为什么会破坏基于统一线性近似的一阶算法解释。
\end{exercise}

\begin{exercise}[Hilbert 梯度桥接题占位]\label{exr:hilbert-gradient-placeholder}
补一题：对例~\ref{ex:hilbert-least-squares} 验证定义~\ref{def:hilbert-gradient}，并说明当 $H_1=\mathbb R^n$、$H_2=\mathbb R^m$ 时，梯度公式如何退化为 Boyd 书中的矩阵表达。
\end{exercise}


\section{对偶空间中的次微分}

\label{sec:5-3}

上一节的可微理论告诉我们：在凸可微点，导数给出局部支撑超平面。然而凸优化里真正有趣的模型往往恰恰出现在不可微处，例如范数、指示函数、绝对值型正则化以及各种约束罚项。此时“唯一导数”不再存在，但支撑超平面仍然可能存在，并且往往不止一个。次微分的作用，就是把这些可能的局部线性支撑统一收集到对偶空间里，从而把不可微问题重新纳入一阶最优性与对偶几何的框架。

\subsection{次微分的定义}

\begin{definition}[次微分]\label{def:subdifferential}
设 $X$ 为 Banach 空间，$f\colon X\to \mathbb R\cup\{+\infty\}$ 为 proper 凸泛函，$x\in \operatorname{dom}f$。定义
\[
\partial f(x)
:=\left\{x^\ast\in X^\ast:
f(y)\ge f(x)+\langle x^\ast,y-x\rangle,\ \forall y\in X\right\}.
\]
集合 $\partial f(x)$ 称为 $f$ 在 $x$ 处的\emph{次微分}；其中元素称为\emph{次梯度}。若 $x\notin \operatorname{dom}f$，则约定 $\partial f(x)=\varnothing$。
\end{definition}

\begin{remark}
定义~\ref{def:subdifferential} 正是第 \ref{sec:4-3} 节定理~\ref{thm:hahn-banach-geometric} 与命题~\ref{prop:supporting-hyperplane} 的局部化版本。区别只在于：这里支撑对象不再是单个凸集，而是凸泛函的上图像
\[
\operatorname{epi}f=\{(x,r)\in X\times\mathbb R:r\ge f(x)\}.
\]
因此次微分可以理解为“在函数图像下方放置的所有支撑超平面之斜率”。
\end{remark}

\begin{proposition}[次梯度最优性条件]\label{prop:subgradient-optimality}
设 $f\colon X\to \mathbb R\cup\{+\infty\}$ 为 proper 凸泛函，$x\in \operatorname{dom}f$。则
\[
0\in \partial f(x)
\]
当且仅当 $x$ 是 $f$ 的全局极小点。
\end{proposition}

\begin{proof}
若 $0\in \partial f(x)$，由定义~\ref{def:subdifferential}，
\[
f(y)\ge f(x),\qquad \forall y\in X,
\]
故 $x$ 为全局极小点。反之，若 $x$ 为全局极小点，则同样有
\[
f(y)\ge f(x)=f(x)+\langle 0,y-x\rangle,\qquad \forall y\in X,
\]
因而 $0\in\partial f(x)$。
\end{proof}

\begin{proposition}[可微点的单值次微分]\label{prop:frechet-subdifferential-singleton}
设 $X$ 为 Banach 空间，$U\subset X$ 为凸开集，$f\colon U\to \mathbb R$ 为凸函数。若 $f$ 在 $x\in U$ 处 Fréchet 可微，则
\[
\partial f(x)=\{Df(x)\}.
\]
在 Hilbert 空间中，这等价于
\[
\partial f(x)=\{\nabla f(x)\},
\]
其中梯度由定义~\ref{def:hilbert-gradient} 给出。
\end{proposition}

\begin{proof}
由上一节的凸可微支撑不等式，$Df(x)$ 满足
\[
f(y)\ge f(x)+Df(x)(y-x),\qquad \forall y\in U.
\]
凸性又保证该不等式可延拓为对所有 $y\in X$ 的次梯度不等式，因此 $Df(x)\in\partial f(x)$。

现取任意 $x^\ast\in \partial f(x)$。由定义~\ref{def:subdifferential}，对任意 $h\in X$ 与充分小的 $t>0$，
\[
f(x+th)\ge f(x)+t\langle x^\ast,h\rangle.
\]
同理，把 $h$ 换成 $-h$ 得
\[
f(x-th)\ge f(x)-t\langle x^\ast,h\rangle.
\]
两式分别除以 $t$ 并令 $t\downarrow 0$，利用 Fréchet 可微得到
\[
Df(x)h\ge \langle x^\ast,h\rangle,\qquad
Df(x)h\le \langle x^\ast,h\rangle.
\]
故对任意 $h\in X$，
\[
Df(x)h=\langle x^\ast,h\rangle.
\]
于是 $x^\ast=Df(x)$。故次微分是单点集。
\end{proof}

\subsection{次微分的非空性}

\begin{theorem}[内点处次微分非空]\label{thm:subdifferential-nonempty-interior}
设 $X$ 为 Banach 空间，$f\colon X\to \mathbb R\cup\{+\infty\}$ 为 proper、凸且下半连续的泛函。若
\[
x\in \operatorname{int}(\operatorname{dom}f),
\]
则
\[
\partial f(x)\neq \varnothing.
\]
\end{theorem}

\begin{proof}
考虑乘积空间 $X\times \mathbb R$，赋予范数
\[
\|(u,r)\|:=\|u\|+|r|.
\]
因为 $f$ 下半连续且凸，其上图像
\[
\operatorname{epi}f:=\{(u,r):r\ge f(u)\}
\]
是闭凸集。点 $(x,f(x))$ 位于其边界上；另一方面，由 $x\in \operatorname{int}(\operatorname{dom}f)$ 可知存在半径 $\rho>0$ 使 $B(x,\rho)\subset \operatorname{dom}f$，因此对任意足够小的 $h$ 与任意 $\eta>0$，点
\[
(x+h,f(x)+1+\eta)
\]
属于 $\operatorname{epi}f$ 的内部，所以该上图像具有非空内部。

于是可对闭凸集 $\operatorname{epi}f$ 在边界点 $(x,f(x))$ 处应用第 \ref{sec:4-3} 节命题~\ref{prop:supporting-hyperplane}。存在非零连续线性泛函 $(x^\ast,s)\in X^\ast\times \mathbb R$ 使
\[
\langle x^\ast,u\rangle + sr
\le
\langle x^\ast,x\rangle + s f(x),
\qquad \forall (u,r)\in \operatorname{epi}f.
\]
对任意 $u\in \operatorname{dom}f$，取 $r=f(u)$，得
\[
\langle x^\ast,u-x\rangle + s(f(u)-f(x))\le 0.
\]
下面说明 $s<0$。若 $s>0$，则对固定的 $u$ 令 $r\to +\infty$，上式不可能始终成立。若 $s=0$，则
\[
\langle x^\ast,u-x\rangle\le 0,\qquad \forall u\in \operatorname{dom}f.
\]
但 $x$ 是定义域内点，可在 $x$ 的任意小邻域中取到 $u=x\pm th$，从而推出 $\langle x^\ast,h\rangle=0$ 对所有 $h\in X$ 成立，即 $x^\ast=0$，与 $(x^\ast,s)\neq 0$ 矛盾。故 $s<0$。

令
\[
\xi^\ast:=-\frac{x^\ast}{s}\in X^\ast.
\]
则上式改写为
\[
f(u)\ge f(x)+\langle \xi^\ast,u-x\rangle,\qquad \forall u\in X.
\]
因此 $\xi^\ast\in \partial f(x)$，从而 $\partial f(x)\neq\varnothing$。
\end{proof}

\subsection{闭图性质}

\begin{proposition}[次微分的闭图性质]\label{prop:subdifferential-closed-graph}
设 $X$ 为 Banach 空间，$f\colon X\to \mathbb R\cup\{+\infty\}$ 为 proper、凸且下半连续的泛函。若网 $x_\alpha\to x$ 在范数拓扑下收敛，网 $x_\alpha^\ast\rightharpoonup^\ast x^\ast$ 在弱$^\ast$拓扑下收敛，且
\[
x_\alpha^\ast\in \partial f(x_\alpha),\qquad \forall \alpha,
\]
则
\[
x^\ast\in \partial f(x).
\]
\end{proposition}

\begin{proof}
由 $x_\alpha^\ast\in \partial f(x_\alpha)$，对任意固定 $y\in X$ 都有
\[
f(y)\ge f(x_\alpha)+\langle x_\alpha^\ast,y-x_\alpha\rangle.
\]
由下半连续性的网刻画，
\[
f(x)\le \liminf_\alpha f(x_\alpha).
\]
另一方面，由 $x_\alpha\to x$ 强收敛以及 $x_\alpha^\ast\rightharpoonup^\ast x^\ast$，对固定的 $y$ 有
\[
\langle x_\alpha^\ast,y-x_\alpha\rangle
\to
\langle x^\ast,y-x\rangle.
\]
因此对上式取上极限得到
\[
f(y)\ge f(x)+\langle x^\ast,y-x\rangle,\qquad \forall y\in X.
\]
由定义~\ref{def:subdifferential}，即 $x^\ast\in \partial f(x)$。
\end{proof}

\subsection{典型例子}

\begin{example}[范数函数的次微分]\label{ex:norm-subdifferential}
设 $X$ 为 Banach 空间，$f(x)=\|x\|$。当 $x\neq 0$ 时，
\[
\partial \|x\|
=
\left\{x^\ast\in X^\ast:\|x^\ast\|=1,\ \langle x^\ast,x\rangle=\|x\|\right\};
\]
而在原点处，
\[
\partial \|0\|=\{x^\ast\in X^\ast:\|x^\ast\|\le 1\}.
\]
把这个例子放在这里，是为了说明不可微并不意味着“一阶信息消失”；相反，它意味着可能存在整片支撑超平面。这里的次梯度集合正是所有“支撑单位球”的对偶向量。
\end{example}

\begin{example}[$\ell_1$ 范数的次梯度]\label{ex:l1-subgradient}
在 $\mathbb R^n$ 上，
\[
f(x)=\|x\|_1=\sum_{i=1}^n |x_i|
\]
满足
\[
\partial \|x\|_1
=
\left\{g\in\mathbb R^n:
g_i=
\begin{cases}
\operatorname{sign}(x_i), & x_i\neq 0,\\
\xi_i\in[-1,1], & x_i=0
\end{cases}
\right\}.
\]
把这个例子放在这里，是为了与 Boyd 书中最常见的非光滑例子对齐，也说明稀疏性结构为什么会对应一个多值的一阶条件。
\end{example}

\begin{example}[指示函数与法锥原型]\label{ex:indicator-normal-cone-preview}
设 $C\subset X$ 为闭凸集，其指示函数定义为
\[
\iota_C(x)=
\begin{cases}
0, & x\in C,\\
+\infty, & x\notin C.
\end{cases}
\]
则 $\partial \iota_C(x)$ 由所有满足
\[
\langle x^\ast,y-x\rangle\le 0,\qquad \forall y\in C
\]
的对偶向量组成。把这个例子放在这里，是为了说明约束并不会在次微分框架外另起炉灶；它们只是通过指示函数把“可行性”编码成了法锥条件。第 \ref{sec:5-4} 节会把这一观察写成正式的演算公式。
\end{example}

\subsection*{有限维对照}

Boyd 第 3 章和第 5 章里把次梯度解释为图像上所有可能的支撑超平面斜率。在 $\mathbb R^n$ 中，这个图像直觉几乎已经足够；但在无穷维里，真正稳固的对象是定义~\ref{def:subdifferential} 中的对偶空间集合，以及定理~\ref{thm:hahn-banach-geometric}、命题~\ref{prop:supporting-hyperplane} 所保证的分离与支撑结构。本节的目的，就是把这些欧氏图像提升为可在 Banach 与 Hilbert 空间中直接调用的一阶语言。

\subsection*{习题（待补）}

\begin{exercise}[次梯度最优性占位]\label{exr:subgradient-optimality-placeholder}
补一题：对一个具体不可微凸函数验证命题~\ref{prop:subgradient-optimality}，并把条件 $0\in\partial f(x)$ 翻译成有限维中的 KKT 风格一阶最优性条件。
\end{exercise}

\begin{exercise}[法锥桥接题占位]\label{exr:normal-cone-placeholder}
补一题：从例~\ref{ex:indicator-normal-cone-preview} 出发，把闭凸集的约束最优性写成指示函数的次微分条件，并解释它与支撑超平面的关系。
\end{exercise}


\section{次微分演算}

\label{sec:5-4}

单个点上的次微分只提供局部信息；真正让它成为建模与算法工具的，是一套可稳定复用的演算规则。和泛函如何分解次梯度？线性观测算子如何把对偶变量从像空间拉回原空间？约束指示函数与法锥如何统一？这些问题在有限维里常被视为“计算技巧”，但在无穷维里它们直接决定第 \ref{chap:06} 章共轭公式、第 \ref{chap:07} 章 KKT 条件以及第 9 至 10 章分裂算法能否写成统一模板。本节的核心接口就是 Moreau-Rockafellar 和公式及其线性拉回版本。

\subsection{和规则与法锥}

\begin{proposition}[次微分和规则的基本包含式]\label{prop:subdifferential-sum-inclusion}
设 $X$ 为 Banach 空间，$f,g\colon X\to \mathbb R\cup\{+\infty\}$ 为 proper 凸泛函。则对任意 $x\in X$，
\[
\partial f(x)+\partial g(x)\subset \partial (f+g)(x).
\]
\end{proposition}

\begin{proof}
任取
\[
x_1^\ast\in \partial f(x),\qquad x_2^\ast\in \partial g(x).
\]
由定义~\ref{def:subdifferential}，对任意 $y\in X$，
\[
f(y)\ge f(x)+\langle x_1^\ast,y-x\rangle,
\qquad
g(y)\ge g(x)+\langle x_2^\ast,y-x\rangle.
\]
两式相加得
\[
(f+g)(y)\ge (f+g)(x)+\langle x_1^\ast+x_2^\ast,y-x\rangle.
\]
因此
\[
x_1^\ast+x_2^\ast\in \partial (f+g)(x),
\]
结论成立。
\end{proof}

\begin{corollary}[指示函数与法锥]\label{cor:normal-cone-indicator}
设 $C\subset X$ 为非空闭凸集，并定义指示函数
\[
\iota_C(x)=
\begin{cases}
0, & x\in C,\\
+\infty, & x\notin C.
\end{cases}
\]
对 $x\in C$，定义法锥
\[
N_C(x):=\{x^\ast\in X^\ast:\langle x^\ast,y-x\rangle\le 0,\ \forall y\in C\}.
\]
则
\[
\partial \iota_C(x)=N_C(x),\qquad x\in C.
\]
\end{corollary}

\begin{proof}
由定义~\ref{def:subdifferential}，
\[
x^\ast\in \partial \iota_C(x)
\]
当且仅当对任意 $y\in X$，
\[
\iota_C(y)\ge \iota_C(x)+\langle x^\ast,y-x\rangle.
\]
当 $y\notin C$ 时左侧为 $+\infty$，不等式自动成立；当 $y\in C$ 时，$\iota_C(y)=\iota_C(x)=0$，故不等式恰化为
\[
\langle x^\ast,y-x\rangle\le 0,\qquad \forall y\in C.
\]
这正是 $x^\ast\in N_C(x)$。
\end{proof}

\subsection{Moreau--Rockafellar 和公式}

\begin{theorem}[Moreau--Rockafellar 和公式]\label{thm:moreau-rockafellar-sum}
设 $X$ 为 Banach 空间，$f,g\colon X\to \mathbb R\cup\{+\infty\}$ 为 proper、凸且下半连续的泛函，$x\in \operatorname{dom}f\cap \operatorname{dom}g$。若 $f$ 或 $g$ 中至少有一个在 $x$ 处连续，则
\[
\partial (f+g)(x)=\partial f(x)+\partial g(x).
\]
\end{theorem}

\begin{proof}
由命题~\ref{prop:subdifferential-sum-inclusion}，只需证明反向包含。设
\[
x^\ast\in \partial (f+g)(x).
\]
若 $f$ 在 $x$ 处连续，则交换 $f,g$ 的角色即可，故不妨设 $g$ 在 $x$ 处连续。定义
\[
F(u):=f(x+u)-f(x)-\langle x^\ast,u\rangle,
\qquad
G(u):=g(x)-g(x+u).
\]
由 $x^\ast\in\partial(f+g)(x)$，对任意 $u\in X$ 都有
\[
F(u)\ge G(u).
\]

考虑集合
\[
C_1:=\{(u,r)\in X\times\mathbb R:F(u)\le r\},
\qquad
C_2:=\{(u,r)\in X\times\mathbb R:G(u)>r\}.
\]
$C_1$ 是闭凸集，而由于 $g$ 在 $x$ 处连续，$G$ 在原点处上半连续，从而 $C_2$ 是开凸集。两者由上式可知不相交。于是由第 \ref{sec:4-3} 节定理~\ref{thm:strict-separation} 的分离思想，存在非零连续线性泛函 $(u^\ast,s)\in X^\ast\times\mathbb R$ 使
\[
\langle u^\ast,u\rangle + sr_1 \ge \langle u^\ast,v\rangle + sr_2,
\qquad
\forall (u,r_1)\in C_1,\ \forall (v,r_2)\in C_2.
\]
取 $u=v=0$，并利用 $(0,r_1)\in C_1$ 对所有 $r_1\ge 0$ 成立、$(0,r_2)\in C_2$ 对所有 $r_2<0$ 成立，可知 $s>0$。再让 $(u,r_1)=(u,F(u))\in C_1$，并令 $r_2\uparrow G(u)$，得到
\[
\langle u^\ast,u\rangle + sF(u)\ge 0,\qquad
\langle u^\ast,u\rangle + sG(u)\le 0.
\]
于是对任意 $u\in X$，
\[
f(x+u)\ge f(x)+\left\langle x^\ast-\frac{u^\ast}{s},u\right\rangle,
\qquad
g(x+u)\ge g(x)+\left\langle \frac{u^\ast}{s},u\right\rangle.
\]
因此
\[
x^\ast-\frac{u^\ast}{s}\in \partial f(x),
\qquad
\frac{u^\ast}{s}\in \partial g(x),
\]
从而
\[
x^\ast\in \partial f(x)+\partial g(x).
\]
反向包含得证。
\end{proof}

\begin{remark}
定理~\ref{thm:moreau-rockafellar-sum} 中的连续性假设是一种典型资格条件。它的作用不是为了“让证明顺手”，而是为了排除病态分离失败。后文在第 \ref{sec:6-4} 节 Fenchel--Rockafellar 对偶和第 \ref{sec:7-3} 节 Slater 条件中，会看到完全相同的逻辑再次出现。
\end{remark}

\subsection{线性映射下的拉回}

\begin{proposition}[线性映射下的拉回与链式法则]\label{prop:subdifferential-linear-pullback}
设 $X,Y$ 为 Banach 空间，$A\colon X\to Y$ 为有界线性算子，$g\colon Y\to \mathbb R\cup\{+\infty\}$ 为 proper 凸泛函。则对任意 $x\in X$，
\[
A^\ast \partial g(Ax)\subset \partial (g\circ A)(x),
\]
其中伴随算子 $A^\ast$ 由第 \ref{sec:3-3} 节定义~\ref{def:adjoint-operator} 给出。
\end{proposition}

\begin{proof}
任取 $y^\ast\in \partial g(Ax)$。则对任意 $z\in X$，
\[
g(Az)\ge g(Ax)+\langle y^\ast,Az-Ax\rangle.
\]
利用伴随算子定义，
\[
\langle y^\ast,Az-Ax\rangle=\langle A^\ast y^\ast,z-x\rangle.
\]
故
\[
(g\circ A)(z)\ge (g\circ A)(x)+\langle A^\ast y^\ast,z-x\rangle,\qquad \forall z\in X.
\]
由定义~\ref{def:subdifferential}，这说明
\[
A^\ast y^\ast\in \partial(g\circ A)(x).
\]
\end{proof}

\begin{remark}
命题~\ref{prop:subdifferential-linear-pullback} 是本章最常用的链式法则版本。它说明观测空间中的次梯度会先在像空间里计算，再通过伴随算子拉回原空间。第 \ref{sec:7-2} 节的算子约束、第 \ref{sec:10-1} 节的前向--后向分裂，以及许多正则化问题中的“梯度项加观测项”都依赖这一步。
\end{remark}

\subsection{典型例子}

\begin{example}[和泛函 $f+g$ 的次微分分解]\label{ex:sum-rule-splitting}
设 $H$ 为 Hilbert 空间，$f$ 为光滑凸泛函，$g$ 为 proper 下半连续凸泛函。若 $g$ 在某点连续，则由定理~\ref{thm:moreau-rockafellar-sum}，
\[
\partial(f+g)(x)=\{\nabla f(x)\}+\partial g(x).
\]
把这个例子放在这里，是为了直接通向第 \ref{sec:10-1} 节的前向--后向分裂：一个部分用梯度处理，另一个部分用邻近映射处理。
\end{example}

\begin{example}[线性观测下的正则化]\label{ex:linear-observation-regularization}
设 $A\colon X\to Y$ 为有界线性算子，$\varphi\colon Y\to \mathbb R\cup\{+\infty\}$ 为 proper 凸泛函，$\psi\colon X\to \mathbb R\cup\{+\infty\}$ 为正则化项。考虑
\[
F(x)=\varphi(Ax)+\psi(x).
\]
则由命题~\ref{prop:subdifferential-linear-pullback} 与定理~\ref{thm:moreau-rockafellar-sum}，
\[
\partial F(x)\supset A^\ast \partial \varphi(Ax)+\partial \psi(x),
\]
在标准资格条件下可取等号。把这个例子放在这里，是为了统一逆问题、机器学习和约束优化里的“数据项 + 正则项”结构。
\end{example}

\begin{example}[有限维复合凸函数]\label{ex:finite-dimensional-composite}
当 $X=\mathbb R^n$、$Y=\mathbb R^m$ 且 $A$ 是矩阵时，命题~\ref{prop:subdifferential-linear-pullback} 退化为
\[
\partial (\varphi\circ A)(x)\supset A^\top \partial \varphi(Ax).
\]
若 $\varphi$ 可微，则进一步得到
\[
\nabla (\varphi\circ A)(x)=A^\top \nabla \varphi(Ax).
\]
把这个例子放在这里，是为了把 Boyd 中关于复合目标和线性约束的矩阵链式规则放回无穷维的统一框架。
\end{example}

\subsection*{有限维对照}

Boyd 第 3 章和第 6 章里关于复合目标、约束消元和拉格朗日处理的许多规则，在欧氏空间里通常写成“把几个梯度和次梯度加起来，再乘一个转置矩阵”。本节说明，这些公式并不是矩阵世界的偶然巧合，而是命题~\ref{prop:subdifferential-sum-inclusion}、定理~\ref{thm:moreau-rockafellar-sum} 和命题~\ref{prop:subdifferential-linear-pullback} 在有限维 Hilbert 空间中的坍缩。

\subsection*{习题（待补）}

\begin{exercise}[Moreau--Rockafellar 公式占位]\label{exr:moreau-rockafellar-placeholder}
补一题：针对一个具体的“光滑项 + 非光滑项”模型，使用定理~\ref{thm:moreau-rockafellar-sum} 写出其最优性条件，并说明资格条件在哪一步被使用。
\end{exercise}

\begin{exercise}[线性拉回桥接题占位]\label{exr:linear-pullback-placeholder}
补一题：对例~\ref{ex:linear-observation-regularization} 验证命题~\ref{prop:subdifferential-linear-pullback}，并把所得结论退回到矩阵情形，与 Boyd 书中的转置矩阵公式逐项对照。
\end{exercise}

